\section{需求建模}
\begin{frame}{业务需求与用户需求}
\begin{block}{业务流程}
准备测试环境 $\rightarrow$ 录入场景 $\rightarrow$ 执行测试 $\rightarrow$ 核对证据 $\rightarrow$ 导出结论与复现材料
\end{block}
\begin{tabularx}{\textwidth}{lX}
\toprule
用户角色 & 主要需求与用例 \\
\midrule
测试人员 & 输入业务场景，启动或取消测试，查看每个检查点的实际结果。\\
开发人员 & 根据失败动作、页面截图与 Trace 复现问题，区分应用缺陷与执行异常。\\
课程评审者 & 复核实验条件、运行证据和评测方法，查看团队协作与交付记录。\\
\bottomrule
\end{tabularx}
\vfill
\small 这些是使用角色；当前系统采用本机使用方式，不对应多用户权限系统。
\end{frame}

\begin{frame}{功能需求与验收依据}
\small
\begin{tabularx}{\textwidth}{lXX}
\toprule
编号 & 功能需求 & 验收依据 \\
\midrule
F1 场景管理 & 新建、编辑、导入场景，关联被测应用 & 原始需求完整保存，历史运行不随编辑改变。\\
F2 自主执行 & 解析、规划、观察、执行与有限恢复 & 真实浏览器可完成代表场景，动作有记录。\\
F3 检查点验证 & 按证据输出三类结论 & 报告列明预期、实际观察与证据，异常不直接判为缺陷。\\
F4 运行跟踪 & 时间线、步数、取消与刷新补读 & 取消后保留记录，刷新后可继续查看已发生事件。\\
F5 报告管理 & 查询历史，下载报告和 Trace & 报告可追溯到场景、环境及运行过程。\\
\bottomrule
\end{tabularx}
\end{frame}

\begin{frame}{非功能需求与课程约束}
\small
\begin{tabularx}{\textwidth}{lX}
\toprule
维度 & 设计约束与检查方式 \\
\midrule
资源与终止 & 浏览器操作最多 35 次，初始导航及失败尝试计数；验证第 36 次操作被拒绝，超时后可收尾。\\
可靠性 & 中断、无效动作、模型超时应保留报告；证据不足输出无法判定。\\
可追溯性 & 记录应用版本、模型信息、时间、动作及证据；固定依赖与被测版本。\\
安全与隐私 & 密钥保存在本机环境配置，截图遮罩密码框，运行证据留在本地。\\
可维护性 & UI 与 CLI 共用核心服务，采用类型约束、自动化检查和模块边界。\\
可用性 & 中文界面展示运行状态及原因，支持新成员按文档初始化并执行示例。\\
\bottomrule
\end{tabularx}
\vfill
\footnotesize 课程约束沿用原报告：指定 DeepSeek 模型与三个固定版本应用；正式评测前复核模型映射，并取得教师调试场景。
\end{frame}
